\documentclass[11pt,a4paper]{article}
\usepackage[utf8]{inputenc}
\usepackage[T1]{fontenc}
\IfFileExists{lmodern.sty}{\usepackage{lmodern}}{\IfFileExists{mathptmx.sty}{\usepackage{mathptmx}}{}}
\usepackage[ngerman]{babel}
\usepackage[a4paper,margin=27mm]{geometry}
\usepackage{amsmath,amssymb}
\usepackage{booktabs}
\usepackage{array}
\renewcommand{\floatpagefraction}{0.85}
\renewcommand{\textfraction}{0.1}
\IfFileExists{enumitem.sty}{\usepackage{enumitem}\setlist[itemize]{leftmargin=1.5em,itemsep=2pt}}{}
\usepackage{tikz}
\usetikzlibrary{arrows.meta,positioning,calc,shapes.misc}
\IfFileExists{microtype.sty}{\usepackage{microtype}}{}
\usepackage[hidelinks]{hyperref}
\setlength{\parskip}{0.35em}
\setlength{\parindent}{1.2em}
\newcommand{\SO}{Subjekt-Objekt}
\newcommand{\Werke}[1]{(W~2, S.~#1)}
\newcolumntype{R}[1]{>{\raggedright\arraybackslash}p{#1}}

\title{Zwei Rekursionen\\[0.4em]
\large Subjekt-Objekt und Identität und Nichtidentität in Hegels Differenzschrift (1801)\\ und die Einheit beider Unterscheidungen}
\author{}
\date{}

\begin{document}
\maketitle

\section{Ort und Anlage}

Die Differenzschrift gliedert sich in eine Vorerinnerung, das Kapitel über die mancherlei Formen des gegenwärtigen Philosophierens, die Darstellung des Fichteschen Systems, die Vergleichung des Schellingschen Prinzips mit dem Fichteschen und den Anhang über Reinhold. Aus diesem Aufbau werden zwei Rekursionen herausgehoben, und dann wird gefragt, ob sie eine sind.

Die erste Rekursion ist die des \SO{}. Das Absolute ist \SO{}; das Subjekt ist \SO{}; das Objekt ist \SO{}; und jeder Teil jedes der beiden ist es wieder. Sie wird in der Vergleichung entwickelt \Werke{94--101} und bis zum Indifferenzpunkt der beiden Wissenschaften geführt \Werke{111~ff.}; vorbereitet ist sie durch das Urteil der Vorerinnerung, das Prinzip Fichtes erweise sich als ein bloß subjektives \SO{} \Werke{10~f.}.

Die zweite Rekursion ist die von Identität und Nichtidentität. Das Absolute ist die Identität der Identität und der Nichtidentität \Werke{96}. Dass dies eine Rekursion ist und nicht bloß eine Formel, zeigt das Kapitel über das Prinzip einer Philosophie in der Form eines absoluten Grundsatzes \Werke{35--41}: Der Satz der Identität enthält in sich die Differenz, der Satz der Nichtidentität enthält in sich die Identität, und beide bedingen einander. Jedes der beiden ist wieder beides, wie im ersten Fall jedes der beiden wieder \SO{} ist.

Die Erweiterung, die den dritten Teil ausmacht, stellt die Frage nach der Einheit der beiden Unterscheidungen. Wenn das Absolute die Identität von Subjekt und Objekt ist und zugleich die Identität von Identität und Nichtidentität, dann ist zu fragen, welche der beiden Unterscheidungen an der Wurzel steht und welche als Argument; ob an der Wurzel Identität steht oder Nichtidentität oder beides; und ob die eine Unterscheidung in den Räumen der anderen wieder auftritt, also Identität und Nichtidentität auf Seiten des Absoluten, des Subjekts und des Objekts, und Subjekt und Objekt innerhalb der Identität und innerhalb der Nichtidentität. Die Antwort, die der Text hergibt, ist, dass alle diese Übergänge stattfinden und dass jede Unterscheidung beim Eintritt in den fremden Raum ihre Form wandelt, ohne ihren Gehalt zu verlieren. Darin besteht ihre Einheit.

Zitiert wird nach Band~2 der Theorie-Werkausgabe (Suhrkamp 1986), abgekürzt W~2, mit der dort eingedruckten Paginierung.

%=====================================================================
\section{Erste Rekursion: \SO{} in Absolutem, Subjekt und Objekt}

\subsection{Der Ausgangspunkt: Entzweiung}

Der Ausgangspunkt ist ein Riss. \glqq Entzweiung ist der Quell des Bedürfnisses der Philosophie\grqq{} \Werke{20}. Die Bildung hat das, was Erscheinung des Absoluten ist, vom Absoluten isoliert und als Selbständiges fixiert. Die feste Gestalt dieser Entzweiung ist die Entgegensetzung von Subjekt und Objekt. Die Aufgabe der Philosophie ist, sie aufzuheben, das Absolute für das Bewusstsein zu konstruieren. Ob das gelingt, entscheidet sich daran, wie die beiden Entgegengesetzten in das Absolute gesetzt werden.

\subsection{Das Urteil über Fichte}

Fichtes Prinzip ist die transzendentale Anschauung, Ich~=~Ich: das Selbstbewusstsein, das sich selbst zum Gegenstand hat. In ihm sind das Subjekt, das setzt, und das Objekt, das gesetzt wird, dasselbe. Das ist der Sinn des Wortes \SO{}: eine Identität, in der Subjekt und Objekt nicht nachträglich zusammengebracht, sondern von Anfang an eins sind. Aber sobald die Spekulation sich zum System bildet, verlässt sie ihr Prinzip. Das Deduzierte erhält die Form einer Bedingung des reinen Bewusstseins; Ich~=~Ich wird durch eine objektive Unendlichkeit bedingt, einen Progress ins Unendliche. \glqq Das Prinzip, das \SO{} erweist sich als ein subjektives \SO{}\grqq{} \Werke{10~f.}. Die Identität von Subjekt und Objekt ist nur unter der Form des Subjekts gesetzt; ihr steht das Objekt als das bloß Gesetzte gegenüber.

\glqq Weil Ich subjektives \SO{} ist, so bleibt ihm eine Seite, von welcher ihm ein Objekt absolut entgegengesetzt ist, von welcher es durch dasselbe bedingt ist\grqq{} \Werke{68--73}. Vom Standpunkt des Ich erscheint das Objekt \glqq als ein nicht Sich-selbst-Bestimmendes, als ein absolut Bestimmtes\grqq{} \Werke{96}. Darum verwandelt sich der Grundsatz: Ich~=~Ich wird zu \glqq Ich soll gleich Ich sein\grqq{} \Werke{97}. Das Absolute des Systems ist ein Sollen.

\subsection{Die Forderung und die Rekursion}

Die Vergleichung setzt mit der Diagnose ein, bei Fichte trete Subjekt~=~Objekt aus der Identität heraus und vermöge sich zu ihr nicht wiederherzustellen, weil das Differente ins Kausalitätsverhältnis versetzt wurde. Dann die Forderung:

\begin{quote}
\glqq Daß absolute Identität das Prinzip eines ganzen Systems sei, dazu ist notwendig, daß das Subjekt und Objekt beide als \SO{} gesetzt werden. Die Identität hat sich im Fichteschen System nur zu einem subjektiven \SO{} konstituiert. Dies bedarf zu seiner Ergänzung eines objektiven \SO{}s, so daß das Absolute sich in jedem der beiden darstellt, vollständig sich nur in beiden zusammen findet, als höchste Synthese in der Vernichtung beider, insofern sie entgegengesetzt sind, als ihr absoluter Indifferenzpunkt beide in sich schließt, beide gebiert und sich aus beiden gebiert.\grqq{} \Werke{94}
\end{quote}

Drei Schritte sind zu lesen. Subjekt und Objekt sind beide \SO{}. Das Absolute stellt sich in jedem dar und ist vollständig nur beide zusammen. Und es ist Indifferenzpunkt, der beide gebiert und sich aus beiden gebiert: Die Richtung der Erzeugung ist doppelt.

Zwei Seiten später wird begründet und die Rekursion ausgesprochen. Die Philosophie muss dem Trennen sein Recht widerfahren lassen, kann die Getrennten aber nicht setzen, ohne sie im Absoluten zu setzen; beide müssen mit gleichem Recht und gleicher Notwendigkeit gesetzt werden, denn würde nur eines auf das Absolute bezogen, wäre die Vereinigung unmöglich \Werke{96}. Dann:

\begin{quote}
\glqq Es müssen daher beide in das Absolute oder das Absolute in beiden Formen gesetzt werden und zugleich beide als Getrennte bestehen; das Subjekt ist hiermit subjektives \SO{}, das Objekt objektives \SO{}. Und weil nunmehr, da eine Zweiheit gesetzt ist, jedes der Entgegengesetzten ein sich selbst Entgegengesetztes ist und die Teilung ins Unendliche geht, so ist jeder Teil des Subjekts und jeder Teil des Objekts selbst im Absoluten, eine Identität des Subjekts und Objekts, -- jedes Erkennen eine Wahrheit, so wie jeder Staub eine Organisation. Nur indem das Objekt selbst ein \SO{} ist, ist Ich~=~Ich das Absolute.\grqq{} \Werke{96~f.}
\end{quote}

Die Zweiheit von Subjekt und Objekt kehrt in jedem der beiden wieder und in jedem Teil jedes der beiden. Die Teilung geht ins Unendliche, ohne dass die Identität verloren ginge. Das Kleinste ist nicht Rest der Teilung, sondern wieder das Ganze.

\begin{figure}[!htb]
\centering
\resizebox{\linewidth}{!}{%
\begin{tikzpicture}[
  every node/.style={align=center, font=\small},
  box/.style={draw, rounded corners=2pt, inner sep=5pt, minimum width=3.2cm},
  leaf/.style={draw, rounded corners=2pt, inner sep=3pt, font=\footnotesize},
  edge/.style={-{Latex[length=2mm]}, thick},
  back/.style={-{Latex[length=2mm]}, thick, dashed, gray!70!black},
  node distance=9mm and 7mm]
\node[box, fill=gray!12] (A) {\textbf{Absolutes}\\ \SO{}\\ Indifferenzpunkt: Identität \emph{und} Totalität};
\node[box, below left=14mm and 6mm of A] (S) {\textbf{Subjekt}\\ subjektives \SO{}\\ Ich = Ich, Intelligenz\\ \footnotesize Subjekt Substanz, Objekt Akzidens};
\node[box, below right=14mm and 6mm of A] (O) {\textbf{Objekt}\\ objektives \SO{}\\ Natur\\ \footnotesize Objekt Substanz, Subjekt Akzidens};
\node[leaf, below left=9mm and -6mm of S] (SS) {Subjekt\\ (im Subjekt)};
\node[leaf, below right=9mm and -6mm of S] (SO) {Objekt\\ (im Subjekt)};
\node[leaf, below left=9mm and -6mm of O] (OS) {Subjekt\\ (im Objekt)};
\node[leaf, below right=9mm and -6mm of O] (OO) {Objekt\\ (im Objekt)};
\draw[edge] (A) -- (S); \draw[edge] (A) -- (O);
\draw[edge] (S) -- (SS); \draw[edge] (S) -- (SO);
\draw[edge] (O) -- (OS); \draw[edge] (O) -- (OO);
\foreach \n in {SS,SO,OS,OO} { \draw[dotted, thick] (\n.south) -- ++(0,-6mm); }
\node[below=7mm of $(SO.south)!0.5!(OS.south)$, font=\footnotesize\itshape] {Teilung ins Unendliche: jeder Teil wieder \SO{}, \glqq jeder Staub eine Organisation\grqq};
\draw[back] (S.north) to[bend left=25] node[left, font=\scriptsize, xshift=-1mm] {gebiert sich\\ aus beiden} (A.west);
\draw[back] (O.north) to[bend right=25] node[right, font=\scriptsize, xshift=1mm] {gebiert sich\\ aus beiden} (A.east);
\end{tikzpicture}}
\caption{Die erste Rekursion nach W~2, S.~94 und 96~f. Abwärts: Das Absolute \glqq gebiert beide\grqq{}, jede Seite ist wieder \SO{}, die Teilung geht ins Unendliche. Aufwärts (gestrichelt): Das Absolute \glqq gebiert sich aus beiden\grqq{}. Die Unterschrift der Seiten nennt das Substantialitätsverhältnis nach S.~101.}
\label{fig:baum}
\end{figure}

\subsection{Reelle Entgegensetzung, die beiden Wissenschaften, der Indifferenzpunkt}

Warum muss die Zweiheit so tief hinabreichen? Weil sonst die Entgegensetzung nicht real wäre. \glqq Indem das Subjekt sowohl als das Objekt ein \SO{} sind, ist die Entgegensetzung des Subjekts und Objekts eine reelle Entgegensetzung; denn beide sind im Absoluten gesetzt und haben dadurch Realität\grqq{} \Werke{97}. Eine ideelle Entgegensetzung ist Werk der Reflexion, die von der absoluten Identität abstrahiert; eine reelle Entgegensetzung ist Werk der Vernunft, die beide Seiten \glqq nicht bloß in der Form des Erkennens, sondern auch in der Form des Seins\grqq{} setzt \Werke{98}. Gerade dadurch, dass beide Seiten dieselbe Identität in sich tragen, wird ihr Gegensatz wirklich. \glqq Hierin besteht allein die wahre Identität, daß beide ein \SO{} sind, und zugleich die wahre Entgegensetzung, deren sie fähig sind\grqq{} \Werke{99}.

Daraus folgen die beiden Wissenschaften. Die Transzendentalphilosophie ist die Wissenschaft vom subjektiven \SO{}, die Naturphilosophie die vom objektiven \SO{} \Werke{100~f.}. Natur und Selbstbewusstsein sind an sich so, wie sie in ihrer Wissenschaft gesetzt werden, \glqq weil es die Vernunft ist, die sie setzt, und die Vernunft setzt sie als \SO{}, also als das Absolute\grqq{}; sie setzt sie so, \glqq weil sie es selbst ist, die sich als Natur und als Intelligenz produziert und sich in ihnen erkennt\grqq{} \Werke{100}. In beiden Wissenschaften stehen das Subjektive und das Objektive im Substantialitätsverhältnis: in der einen die Intelligenz Substanz und die Natur Akzidens, in der anderen umgekehrt \Werke{101}. Der Unterschied von Subjekt und Objekt ist kein Unterschied des Gehalts, denn beide sind \SO{}, sondern ein Unterschied dessen, welche Seite Substanz ist; darum ist die reelle Entgegensetzung \glqq nur quantitativ\grqq{}, und das Absolute, das sich aus ihr rekonstruiert, \glqq kein Quantum, sondern Totalität\grqq{} \Werke{99}.

Beide Wissenschaften sind relative Totalitäten; als solche \glqq streben sie nach dem Indifferenzpunkt; als Identität und als relative Totalität liegt er überall in ihnen selbst, als absolute Totalität außer ihnen\grqq{} \Werke{111}. Sie hängen \glqq als Pole der Indifferenz in dieser selbst zusammen; sie selbst sind die Linien, welche den Pol mit dem Mittelpunkt verknüpfen\grqq{}. Und: \glqq dieser Mittelpunkt ist selbst ein gedoppelter, einmal Identität, das andere Mal Totalität\grqq{} \Werke{111}.

Der Satz vom gedoppelten Mittelpunkt lässt sich an einer einfachen Gestalt nachprüfen. Man nehme eine offene Strecke: zwei Enden, die nicht dazugehören, und eine Mitte, die dazugehört. Die Enden seien die Pole, die Mitte der Indifferenzpunkt, und jeder Punkt dazwischen habe einen Anteil an beiden Seiten, so dass die Differenz nur quantitativ ist. Teilt man die Strecke in der Mitte, so ist jede Hälfte wieder eine offene Strecke mit zwei Enden und eigener Mitte, und dasselbe gilt für jeden Teil jedes Teils. Jeder Teil ist dem Ganzen der Gestalt nach gleich, ohne verkleinert zu sein. Der Mittelpunkt ist in dieser Figur gedoppelt, genau wie der Satz es verlangt: Jeder Teil hat seine eigene Mitte in seinem Inneren; die Mitte des Ganzen aber liegt auf keiner Stufe der Teilung im Inneren eines Teils, sie ist stets nur Rand von genau zwei Teilen (Abbildung~\ref{fig:strecke}). Als Identität liegt der Indifferenzpunkt überall in den Teilen, als absolute Totalität außer ihnen. Die Figur zeigt zugleich Fichtes Fehler: Der Pol allein ist ein Punkt ohne Inneres, ohne Mitte, ohne Teile. Wer ihn als das Ganze nimmt, nimmt einen Rand für den Raum, den er begrenzt. Das subjektive \SO{} hingegen ist die Hälfte, die diesen Pol als Rand hat; sie hat ihre eigene Mitte und ist darum selbst eine Identität von Subjekt und Objekt.

\begin{figure}[!htb]
\centering
\resizebox{\linewidth}{!}{%
\begin{tikzpicture}[x=5.2cm, y=1cm, every node/.style={font=\small, align=center}]
\draw[thick] (-1,0) -- (1,0);
\draw[fill=white, thick] (-1,0) circle (2.2pt) node[below=3pt] {Pol: Subjekt};
\draw[fill=white, thick] (1,0) circle (2.2pt) node[below=3pt] {Pol: Objekt};
\fill (0,0) circle (2.6pt) node[above=3pt] {Indifferenzpunkt des Ganzen};
\node[left=6mm of {(-1,0)}, font=\scriptsize] {Ganzes};
\foreach \a/\b in {-1/0, 0/1} {
  \draw[thick] (\a,-1.3) -- (\b,-1.3);
  \draw[fill=white, thick] (\a,-1.3) circle (2.2pt);
  \draw[fill=white, thick] (\b,-1.3) circle (2.2pt);
  \fill[gray!60!black] ({(\a+\b)/2},-1.3) circle (2.2pt);
}
\node[above=2pt, font=\scriptsize] at (-0.5,-1.3) {subjektives \SO{}\\ eigene Mitte};
\node[above=2pt, font=\scriptsize] at (0.5,-1.3) {objektives \SO{}\\ eigene Mitte};
\node[left=6mm of {(-1,-1.3)}, font=\scriptsize] {erste Teilung};
\foreach \a/\b in {-1/-0.5, -0.5/0, 0/0.5, 0.5/1} {
  \draw[thick] (\a,-2.4) -- (\b,-2.4);
  \draw[fill=white, thick] (\a,-2.4) circle (2.2pt);
  \draw[fill=white, thick] (\b,-2.4) circle (2.2pt);
  \fill[gray!60!black] ({(\a+\b)/2},-2.4) circle (2.2pt);
}
\node[left=6mm of {(-1,-2.4)}, font=\scriptsize] {zweite Teilung};
\foreach \a/\b in {-1/-0.75, -0.75/-0.5, -0.5/-0.25, -0.25/0, 0/0.25, 0.25/0.5, 0.5/0.75, 0.75/1} {
  \draw[thick] (\a,-3.3) -- (\b,-3.3);
  \draw[fill=white, thick] (\a,-3.3) circle (1.8pt);
  \draw[fill=white, thick] (\b,-3.3) circle (1.8pt);
  \fill[gray!60!black] ({(\a+\b)/2},-3.3) circle (1.8pt);
}
\node[left=6mm of {(-1,-3.3)}, font=\scriptsize] {dritte Teilung};
\draw[dashed, gray] (0,0.35) -- (0,-3.7);
\node[below, font=\scriptsize\itshape] at (0,-3.75) {die Mitte des Ganzen: auf jeder Stufe nur Rand, nie Inneres};
\end{tikzpicture}}
\caption{Die Teilung ins Unendliche und der gedoppelte Mittelpunkt nach W~2, S.~96~f. und 111. Offene Kreise: Pole, die nicht zum Teil gehören. Graue Punkte: die eigene Mitte jedes Teils, als Identität überall in den Teilen. Schwarzer Punkt und gestrichelte Linie: die Mitte des Ganzen, als absolute Totalität außer allen Teilen.}
\label{fig:strecke}
\end{figure}

\subsection{Die Formeln der ersten Rekursion}

\begin{center}
\renewcommand{\arraystretch}{1.35}
\begin{tabular}{@{}l R{5.6cm} R{7.4cm}@{}}
\toprule
 & Formel & Wortlaut, auf den sie sich stützt \\
\midrule
(A) & Absolutes $=$ \SO{} & \glqq als ihr absoluter Indifferenzpunkt beide in sich schließt\grqq{} (S.~94) \\
(S) & Subjekt $=$ (\SO{})$_{\text{Form: Subjekt}}$ & \glqq das Subjekt ist hiermit subjektives \SO{}\grqq{} (S.~96) \\
(O) & Objekt $=$ (\SO{})$_{\text{Form: Objekt}}$ & \glqq das Objekt objektives \SO{}\grqq{} (S.~96) \\
(R) & jeder Teil von S, jeder Teil von O $=$ \SO{} & \glqq die Teilung ins Unendliche geht\grqq{} (S.~96~f.) \\
\bottomrule
\end{tabular}
\end{center}

Die Formeln (S) und (O) unterscheiden sich nur im Index; ihr Gehalt ist derselbe. Der Zusatz (R) macht aus ihnen eine Struktur: Selbstähnlichkeit ohne Abbruch, und zwar nicht die schlechte Unendlichkeit, in der Ich~=~Ich sich in einem Progress verliert \Werke{10~f.}, sondern die, in der die Identität auf jeder Stufe schon da ist.

Das Wort Zusammensetzung ist dabei zu korrigieren. Hegel sagt nicht, das Absolute bestehe aus Subjekt und Objekt wie ein Ganzes aus Stücken; er sagt, es stelle sich in jedem dar, finde sich vollständig nur in beiden, gebäre beide und gebäre sich aus beiden \Werke{94}. Die Formulierung vom Gebären hat eine Herkunft, die Hegel selbst nennt: Kant stelle die Natur als ein \SO{} dar, indem er das Naturprodukt als Naturzweck betrachte, \glqq zweckmäßig ohne Begriff, notwendig ohne Mechanismus, Begriff und Sein identisch\grqq{} \Werke{103}. Kant bestimmt ein organisiertes Wesen dadurch, dass die Teile nur durch ihre Beziehung auf das Ganze möglich sind und einander und das Ganze wechselseitig hervorbringen; er hält dies für eine Maxime der Urteilskraft, nicht für ein Prinzip der Natur. Hegels Satz vom Gebären ist diese Bestimmung des Organismus, aus dem Als-ob in das Sein versetzt und vom Lebewesen auf das Ganze übertragen. Darum kann jeder Staub eine Organisation sein, und darum ist die Zusammensetzung keine aus Stücken: Stücke bringen einander nicht hervor.

%=====================================================================
\section{Zweite Rekursion: Identität und Nichtidentität}

\subsection{Die beiden Sätze}

Das Kapitel über den absoluten Grundsatz geht von der Forderung aus, die Philosophie solle ihr Prinzip in einem Satz aussprechen. Ein Satz ist aber eine Form des Verstandes; die Vernunft muss trennen und \glqq die Synthese und die Antithese getrennt, in zwei Sätzen, in einem die Identität, im andern die Entzweiung ausdrücken\grqq{} \Werke{37}.

Der erste Satz ist $A = A$. In ihm \glqq wird reflektiert auf das Bezogensein, und dies Beziehen, dies Einssein, die Gleichheit ist in dieser reinen Identität enthalten; es wird von aller Ungleichheit abstrahiert\grqq{} \Werke{37}. Und: $A = A$, \glqq Ausdruck des absoluten Denkens oder der Vernunft\grqq{}, hat \glqq für die formale, in verständigen Sätzen sprechende Reflexion nur die Bedeutung der Verstandesidentität, der reinen Einheit\grqq{} \Werke{37}. Dasselbe Zeichen hat zwei Bedeutungen.

Der zweite Satz ist $A = B$. Die Vernunft postuliert \glqq auch das Setzen desjenigen, wovon in der reinen Gleichheit abstrahiert wurde, das Setzen des Entgegengesetzten, der Ungleichheit; das eine $A$ ist Subjekt, das andere Objekt, und der Ausdruck für ihre Differenz ist $A$ nicht $= A$, oder $A = B$\grqq{} \Werke{37}. In ihm ist die Nichtidentität gesetzt, \glqq die reine Form des Nichtdenkens\grqq{}. Er ist der Satz des Grundes, herabgezogen zum Satz der Kausalität \Werke{38}.

\subsection{Die Rekursion in den Sätzen}

Dass hier eine Rekursion vorliegt und nicht nur ein Paar, zeigen drei Bestimmungen.

Erstens enthält jeder der beiden Sätze den anderen. Die Nichtidentität kann nur gesetzt werden, weil auch das Nichtdenken gedacht wird; das Beziehen, das Gleichheitszeichen, ist auch im zweiten Satz, \glqq aber nur subjektiv, d.\,h. nur insofern das Nichtdenken durchs Denken gesetzt ist\grqq{} \Werke{38}. Umgekehrt enthält der erste Satz die Differenz: $A = A$, als Beziehung beider Sätze genommen, \glqq enthält die Differenz des $A$ als Subjekts und $A$ als Objekts zugleich mit der Identität, so wie $A = B$ die Identität des $A$ und $B$ mit der Differenz beider\grqq{} \Werke{39}. Die Identität enthält Identität und Nichtidentität; die Nichtidentität enthält Identität und Nichtidentität. Jedes der beiden ist wieder beides, unter der Form des einen. Das ist wörtlich dieselbe Struktur wie im ersten Teil: Das Subjekt ist \SO{} unter der Form des Subjekts, das Objekt \SO{} unter der Form des Objekts. In $A = A$ ist die Identität die Form und die Differenz der Gehalt; in $A = B$ ist die Differenz die Form und die Identität der Gehalt. Das ist das Substantialitätsverhältnis der zweiten Rekursion.

Zweitens bedingen beide einander. \glqq Dieser zweite Satz ist so unbedingt als der erste und insofern Bedingung des ersten, so wie der erste Bedingung des zweiten Satzes ist\grqq{} \Werke{38}. Der erste besteht nur durch die Abstraktion von der Ungleichheit, die der zweite enthält; der zweite braucht, um ein Satz zu sein, das Beziehen, das der erste ausspricht. Das ist das Gegenstück zum Satz, das Absolute gebäre beide und gebäre sich aus beiden: Die beiden Sätze bringen einander hervor.

Drittens ist ihre Beziehung selbst wieder ein Satz, und zwar von derselben Gestalt. \glqq Beide Sätze sind Sätze des Widerspruchs, nur im verkehrten Sinne. Der erste, der der Identität, sagt aus, daß der Widerspruch $= 0$ ist; der zweite, insofern er auf den ersten bezogen wird, daß der Widerspruch ebenso notwendig ist als der Nichtwiderspruch\grqq{} \Werke{38~f.}. Ihre Beziehung ist \glqq der Ausdruck der Antinomie\grqq{}, und als Antinomie, \glqq als Ausdruck der absoluten Identität\grqq{}, ist es gleichgültig, ob $A = B$ oder $A = A$ gesetzt wird \Werke{39}. Die Antinomie ist der \glqq sich selbst aufhebende Widerspruch\grqq{}, der höchste formale Ausdruck des Wissens \Werke{39}; von der Seite der Reflexion angesehen \glqq erscheint die absolute Identität in Synthesen Entgegengesetzter, also in Antinomien\grqq{} \Werke{41}.

\subsection{Fichtes Entscheidung und die Formel}

Fichtes System ist ein bestimmter Umgang mit den beiden Sätzen: \glqq $A = A$ und $A = B$ bleiben beide unbedingt; es soll nur $A = A$ gelten; d.\,h. aber, ihre Identität ist nicht in ihrer wahren Synthese, die kein bloßes Sollen ist, dargestellt\grqq{} \Werke{50}. Die Totalität der Vernunft führt den zweiten Satz herbei, der ein Nicht-Ich setzt; die Synthese wird postuliert, aber \glqq der eine Satz soll bestehen, der eine höher an Rang sein als der andere\grqq{}. Das Sollen ist die Gestalt der nicht vollzogenen Identität der beiden Sätze.

In der Vergleichung kehrt die Zweiheit der Sätze als Zweiheit von Trennung und Identität wieder. Die Philosophie muss dem Trennen sein Recht widerfahren lassen; setzt sie es aber der Identität gleich absolut, so hat sie es nur bedingt gesetzt, wie eine Identität, die durch Vernichten der Entgegengesetzten bedingt ist, nur relativ ist. \glqq Das Absolute selbst aber ist darum die Identität der Identität und der Nichtidentität; Entgegensetzen und Einssein ist zugleich in ihm\grqq{} \Werke{96}. Und zwei Seiten später: Die reelle Entgegensetzung ist Werk der Vernunft, \glqq welche die Entgegengesetzten nicht bloß in der Form des Erkennens, sondern auch in der Form des Seins, Identität und Nichtidentität identisch setzt\grqq{} \Werke{98}.

\subsection{Die Formel als Rekursion}

Die Formel enthält das Wort Identität zweimal, und die beiden Vorkommen bedeuten nicht dasselbe; das ist die Anwendung dessen, was über $A = A$ gesagt wurde. Die innere Identität ist $A = A$, die Verstandesidentität; die Nichtidentität ist $A = B$; die äußere Identität ist die der Vernunft, die beide zusammenhält. Nun sind aber die innere Identität und die Nichtidentität, wie gezeigt, selbst wieder beides. Also lautet die Formel ausgeschrieben: Das Absolute ist die Identität der (Identität von Identität und Nichtidentität) und der (Nichtidentität von Identität und Nichtidentität), und in jeder Klammer geht es weiter. Die zweite Rekursion hat dieselbe Gestalt wie die erste: eine Wurzel, die das Paar als Einheit ist, zwei Kinder, die je das Paar unter einer Form sind, und so fort ohne Ende.

\begin{figure}[!htb]
\centering
\resizebox{\linewidth}{!}{%
\begin{tikzpicture}[
  every node/.style={align=center, font=\small},
  box/.style={draw, rounded corners=2pt, inner sep=5pt, minimum width=3.4cm},
  leaf/.style={draw, rounded corners=2pt, inner sep=3pt, font=\footnotesize},
  edge/.style={-{Latex[length=2mm]}, thick},
  back/.style={-{Latex[length=2mm]}, thick, dashed, gray!70!black}]
\node[box, fill=gray!12] (R) at (0,0) {\textbf{Absolutes}\\ Identität der Identität und der Nichtidentität\\ \footnotesize \glqq Entgegensetzen und Einssein ist zugleich in ihm\grqq{}};
\node[box] (I) at (-3.9,-2.3) {\textbf{Identität}\\ $A = A$\\ \footnotesize Form: Identität, Gehalt: Differenz\\ \footnotesize (\glqq das eine $A$ Subjekt, das andere Objekt\grqq{})};
\node[box] (N) at (3.9,-2.3) {\textbf{Nichtidentität}\\ $A = B$\\ \footnotesize Form: Differenz, Gehalt: Identität\\ \footnotesize (das $=$ \glqq nur subjektiv\grqq{} enthalten)};
\node[leaf] (II) at (-5.6,-4.4) {Identität\\ (in der Identität)};
\node[leaf] (IN) at (-2.2,-4.4) {Nichtidentität\\ (in der Identität)};
\node[leaf] (NI) at (2.2,-4.4) {Identität\\ (in der Nichtidentität)};
\node[leaf] (NN) at (5.6,-4.4) {Nichtidentität\\ (in der Nichtidentität)};
\draw[edge] (R) -- (I); \draw[edge] (R) -- (N);
\draw[edge] (I) -- (II); \draw[edge] (I) -- (IN); \draw[edge] (N) -- (NI); \draw[edge] (N) -- (NN);
\foreach \n in {II,IN,NI,NN} { \draw[dotted, thick] (\n.south) -- ++(0,-5mm); }
\draw[back, <->, >={Latex[length=2mm]}] (I.east) -- node[above=0.5mm, font=\scriptsize] {jeder Bedingung\\ des anderen} (N.west);
\node[below=6mm of $(IN.south)!0.5!(NI.south)$, font=\footnotesize\itshape] {die Beziehung beider: Antinomie; als Antinomie \glqq gleichgültig, $A = B$ oder $A = A$ zu setzen\grqq{}};
\end{tikzpicture}}
\caption{Die zweite Rekursion nach W~2, S.~37--39 und 96. Jeder der beiden Sätze enthält beide Momente unter seiner Form, beide bedingen einander, und ihre Einheit ist die Wurzel. Die Gestalt ist dieselbe wie in Abbildung~\ref{fig:baum}.}
\label{fig:baumIN}
\end{figure}

Ein Unterschied zur ersten Rekursion ist sofort sichtbar, und er wird im dritten Teil tragend. Die Wurzel der ersten heißt \SO{}, mit beiden Namen; die Wurzel der zweiten heißt Identität, mit dem Namen eines der beiden Kinder. Hegel nennt das Absolute nirgends Nichtidentität, aber er nennt es auch nirgends Subjekt. Die Namensgebung ist im einen Fall paritätisch, im anderen nicht. Ein zweiter Unterschied liegt in der Richtung. Die erste Rekursion wird von oben nach unten entwickelt: Das Absolute stellt sich in beiden dar, die Teilung geht ins Unendliche. Die zweite wird von unten nach oben entwickelt: Die beiden Sätze werden getrennt aufgestellt, ihre wechselseitige Bedingtheit wird gezeigt, ihre Antinomie, und daraus ergibt sich das Absolute als ihre Identität. Die eine Rekursion ist Darstellung, die andere Synthese.

%=====================================================================
\section{Die Einheit der beiden Unterscheidungen}

\subsection{Die Frage}

Zwei Unterscheidungen sind rekursiv: Subjekt und Objekt, Identität und Nichtidentität. Beide haben dieselbe Gestalt. Das Absolute wird einmal als Identität von Subjekt und Objekt bestimmt, einmal als Identität von Identität und Nichtidentität. Damit stehen an der Wurzel zwei Fragen. Welches ist der Operator an der Wurzel: Identität, Nichtidentität oder beides? Und welches ist sein Argument: das Paar Subjekt und Objekt oder das Paar Identität und Nichtidentität? Die Frage lautet also: Identität und/oder Nichtidentität von Identität und/oder Nichtidentität, und von Subjekt und Objekt.

Daran schließt die zweite Frage: Treten die Unterscheidungen in den ihnen fremden Räumen wieder auf? Also Identität und Nichtidentität auf Seiten des Absoluten, des Subjekts und des Objekts; und Subjekt und Objekt innerhalb der Identität und innerhalb der Nichtidentität. Und wenn ja, in welcher Form?

\subsection{Der Operator an der Wurzel}

Man gehe die Kombinationen durch und suche zu jeder den Namen, den Hegel ihr gibt.

Identität von Subjekt und Objekt: Das ist die absolute Identität, Subjekt~=~Objekt, der Indifferenzpunkt \Werke{94}. Hegel nennt sie das Prinzip des ganzen Schellingschen Systems.

Nichtidentität von Subjekt und Objekt: Das ist die Entzweiung \Werke{20}, und in ihrer wahren Gestalt die reelle Entgegensetzung \Werke{97}. Sie ist nicht die Verneinung der Identität, sondern durch sie bedingt: \glqq Die Realität Entgegengesetzter und reelle Entgegensetzung findet allein durch die Identität beider statt\grqq{} \Werke{97}.

Identität von Identität und Nichtidentität: Das ist das Absolute nach der Formel \Werke{96}, und in der Form des Erkennens die transzendentale Anschauung, die die Entgegengesetzten nicht bloß denkt, sondern als eins anschaut.

Nichtidentität von Identität und Nichtidentität: Das ist die Antinomie. Der zweite Satz \glqq widerspricht dem vorigen geradezu\grqq{} \Werke{37}; beide sind Sätze des Widerspruchs im verkehrten Sinne; ihre Beziehung ist der sich selbst aufhebende Widerspruch \Werke{38~f.}.

Alle Kombinationen sind im Text besetzt, und keine ist überflüssig. Sie bilden ein Feld (Abbildung~\ref{fig:matrix}). Was das Feld zeigt, ist zuerst dies: Die Nichtidentität an der Wurzel ist nie das Letzte. Die Entzweiung ist der Quell des Bedürfnisses, nicht die Antwort; die Antinomie ist der höchste Ausdruck der Vernunft durch den Verstand, nicht die Vernunft. In beiden Spalten steht die Nichtidentität als das, was die Identität an der Wurzel voraussetzt und was diese in sich aufbewahrt. Das ist der Sinn des Satzes, Entgegensetzen und Einssein sei zugleich im Absoluten. Der Operator an der Wurzel ist darum nicht Identität oder Nichtidentität, sondern Identität und Nichtidentität zugleich, und Hegel nennt dieses Zugleich mit dem Namen der einen Seite, Identität, weil in der Anschauung, nicht in der Reflexion, das Zugleich als Eines gegeben ist. Das Oder ist die Version des Verstandes: \glqq es soll nur $A = A$ gelten\grqq{}. Das Und ist die Version der Vernunft.

\begin{figure}[!htb]
\centering
\resizebox{0.95\linewidth}{!}{%
\begin{tikzpicture}[
  every node/.style={align=center, font=\small},
  cell/.style={draw, rounded corners=2pt, inner sep=6pt, minimum width=5.2cm, minimum height=2.4cm, text width=4.8cm},
  head/.style={font=\small\bfseries}]
\node[head] at (-3.1,1.9) {Argument: Subjekt und Objekt};
\node[head] at (3.1,1.9) {Argument: Identität und Nichtidentität};
\node[head, text width=1.9cm] at (-7.3,0) {Wurzel:\\ Identität};
\node[head, text width=1.9cm] at (-7.3,-2.9) {Wurzel:\\ Nicht-\\ identität};
\node[cell, fill=gray!12] (a) at (-3.1,0) {absolute Identität\\ Subjekt $=$ Objekt, Indifferenzpunkt\\ \footnotesize S.~94, 111};
\node[cell, fill=gray!12] (b) at (3.1,0) {das Absolute nach der Formel\\ transzendentale Anschauung\\ \footnotesize S.~96, 41};
\node[cell] (c) at (-3.1,-2.9) {Entzweiung; reelle Entgegensetzung\\ \glqq allein durch die Identität beider\grqq{}\\ \footnotesize S.~20, 97};
\node[cell] (d) at (3.1,-2.9) {Antinomie\\ der sich selbst aufhebende Widerspruch\\ \footnotesize S.~37--39};
\draw[-{Latex[length=2mm]}, thick, gray!70!black] (c.north) -- node[right=1mm, font=\scriptsize] {bedingt durch} (a.south);
\draw[-{Latex[length=2mm]}, thick, gray!70!black] (d.north) -- node[right=1mm, font=\scriptsize] {aufgehoben in} (b.south);
\draw[<->, >={Latex[length=2mm]}, thick, gray!70!black] (a.east) -- node[above=0.5mm, font=\scriptsize] {Sein /\\ Erkennen} (b.west);
\draw[<->, >={Latex[length=2mm]}, thick, gray!70!black] (c.east) -- node[above=0.5mm, font=\scriptsize] {Sein /\\ Erkennen} (d.west);
\end{tikzpicture}}
\caption{Der Operator an der Wurzel und sein Argument. Jede Kombination hat im Text einen Namen. Die untere Zeile ist in beiden Spalten das Vorausgesetzte und Aufbewahrte der oberen; die beiden Spalten verhalten sich wie die Form des Seins zur Form des Erkennens (W~2, S.~98).}
\label{fig:matrix}
\end{figure}

Und das Feld zeigt ein Zweites. Die beiden Spalten sind nicht zwei Gegenstände, sondern zwei Formen desselben. Die Identität von Subjekt und Objekt ist das Absolute in der Form des Seins: als Natur und Intelligenz, als die beiden Wissenschaften. Die Identität von Identität und Nichtidentität ist das Absolute in der Form des Erkennens: als die Beziehung der beiden Sätze, als Anschauung. Hegel sagt es an der Stelle, an der er die reelle Entgegensetzung bestimmt: Die Vernunft setzt Identität und Nichtidentität identisch \glqq nicht bloß in der Form des Erkennens, sondern auch in der Form des Seins\grqq{} \Werke{98}. Die Nichtidentität, in der Form des Seins gesetzt, ist die reelle Entgegensetzung von Subjekt und Objekt; in der Form des Erkennens ist sie die Antinomie. Die Unterscheidung Subjekt/Objekt ist also die Unterscheidung Identität/Nichtidentität, sofern sie ist; und diese ist jene, sofern sie gewusst wird.

\subsection{Identität und Nichtidentität auf Seiten des Absoluten}

Nun der Eintritt der einen Unterscheidung in die Räume der anderen. Zuerst in den Raum des Absoluten, der Wurzel der ersten Rekursion.

Das Absolute ist dort als Indifferenzpunkt bestimmt. Aber der Indifferenzpunkt ist \glqq selbst ein gedoppelter, einmal Identität, das andere Mal Totalität\grqq{} \Werke{111}. Das ist die Unterscheidung von Identität und Nichtidentität, eingetreten in den Raum des Absoluten und dort in eine neue Form gewandelt. Die Identität ist die Mitte, der Punkt, an dem die Differenz gleichgültig wird. Die Totalität ist das Ganze, das die Differenz enthält: das Absolute, das \glqq sich aus der quantitativen Differenz rekonstruiert\grqq{} und \glqq kein Quantum, sondern Totalität\grqq{} ist \Werke{99}. Totalität ist der Name der Nichtidentität auf Seiten des Absoluten. Sie ist nicht die Entzweiung, denn die Entzweiung fixiert die Getrennten außerhalb des Absoluten; sie ist die Nichtidentität, die im Absoluten aufbewahrt ist, das Entgegensetzen, das zugleich mit dem Einssein in ihm ist. Die Strecke hat das gezeigt: Die Mitte des Ganzen ist Identität, das Ganze mit seinen beiden Polen ist Totalität, und die Mitte liegt außer allen Teilen, das Ganze umfasst sie alle. Auf Seiten des Absoluten heißt Identität Punkt und Nichtidentität Ganzes.

\subsection{Identität und Nichtidentität auf Seiten des Subjekts}

Im Raum des Subjekts, des subjektiven \SO{}, ist die Unterscheidung schon bei Fichte da, und sie hat dort ihre bestimmte Form. Die Identität ist Ich~=~Ich, das Selbstbewusstsein, das Sich-selbst-Setzen \Werke{96~f.}. Die Nichtidentität ist das Nicht-Ich: \glqq Die Totalität der Vernunft führt den zweiten Satz herbei, der ein Nicht-Ich setzt\grqq{} \Werke{50}. Identität und Nichtidentität heißen auf Seiten des Subjekts Selbstbewusstsein und Bewusstsein: das Wissen, das sich selbst weiß, und das Wissen, das ein Anderes weiß. Der Satz $A = A$ ist im Subjekt Ich~=~Ich; der Satz $A = B$ ist im Subjekt Ich und Nicht-Ich.

Und hier zeigt sich, was Fichtes Fehler in dieser Sprache ist. Er hat die beiden Sätze in den Raum des Subjekts eingelassen, aber ihre Identität dort nicht vollzogen; die Synthese wird postuliert, und \glqq in dieser bleibt die Entgegensetzung\grqq{} \Werke{50}. Das Sollen ist die Gestalt, die die Identität der Identität und der Nichtidentität annimmt, wenn sie nur im Raum des Subjekts gesetzt wird: Ich soll gleich Ich sein. Der Grund ist, dass das Nicht-Ich im Subjekt nur als Schranke, als absolut Bestimmtes gesetzt ist, also nicht selbst \SO{}. Die Nichtidentität auf Seiten des Subjekts kann nur dann in die Identität aufgehoben werden, wenn das, was dem Subjekt nicht identisch ist, seinerseits die Struktur des Subjekts hat. Das heißt: Die zweite Rekursion kann im Raum des Subjekts nur geschlossen werden, wenn die erste Rekursion auch nach außen geht. \glqq Nur indem das Objekt selbst ein \SO{} ist, ist Ich~=~Ich das Absolute\grqq{} \Werke{97}. Dieser Satz ist der Schnittpunkt der beiden Rekursionen.

\subsection{Identität und Nichtidentität auf Seiten des Objekts}

Im Raum des Objekts, der Natur, ist die Unterscheidung an der Stelle bestimmt, an der Hegel Kant bespricht. Kant \glqq stellt die Natur als ein \SO{} dar, indem er das Naturprodukt als Naturzweck betrachtet, zweckmäßig ohne Begriff, notwendig ohne Mechanismus, Begriff und Sein identisch\grqq{} \Werke{103}. Die Identität auf Seiten des Objekts heißt Naturzweck, Organisation: Begriff und Sein identisch, das Ganze, das seine Teile hervorbringt. Die Nichtidentität auf Seiten des Objekts heißt Mechanismus: die Notwendigkeit, in der ein Teil den anderen von außen bestimmt, das Kausalitätsverhältnis, die \glqq Menge von Endlichkeiten\grqq{} \Werke{94}, in die das Fichtesche System zerfällt. Der Satz $A = A$ ist in der Natur der Organismus; der Satz $A = B$ ist in der Natur die Kausalität, der Satz des Grundes, den Hegel ausdrücklich als die herabgezogene Form des zweiten Satzes bezeichnet \Werke{38}.

Auch hier zeigt sich ein Fehler, und er ist das genaue Gegenstück zu Fichtes. Kant lässt die Identität in die Natur eintreten, aber nur als Maxime: Die teleologische Ansicht \glqq soll nur als Maxime unseres eingeschränkten, diskursiv denkenden, menschlichen Verstandes gelten\grqq{}, und so bleibt \glqq die Betrachtungsart ein durchaus Subjektives und die Natur ein rein Objektives\grqq{} \Werke{103}. Das heißt: In dem Augenblick, in dem die Identität auf Seiten des Objekts gesetzt werden soll, bricht die Unterscheidung von Subjekt und Objekt wieder auf, und zwar innerhalb des Objekts. Die Zweckmäßigkeit wird dem Subjekt zugeschrieben, der Mechanismus dem Objekt. Fichte kann die Nichtidentität im Subjekt nicht aufheben, Kant die Identität im Objekt nicht setzen, und beide aus demselben Grund: Die eine Rekursion wird ohne die andere geführt.

\begin{figure}[!htb]
\centering
\resizebox{\linewidth}{!}{%
\begin{tikzpicture}[
  every node/.style={align=center, font=\small},
  sp/.style={draw, rounded corners=3pt, inner sep=0pt, minimum width=4.6cm, minimum height=3.4cm},
  half/.style={inner sep=4pt, text width=4.2cm, align=center, font=\footnotesize}]
\foreach \x/\nm/\name/\idt/\nid/\fehl in {
  -6.8/A/{\textbf{Absolutes}}/{Identität:\\ Indifferenzpunkt, die Mitte\\ \footnotesize S.~111}/{Nichtidentität:\\ Totalität, das Ganze mit seinen Polen\\ \footnotesize S.~99, 111}/{},
  0/S/{\textbf{Subjekt}}/{Identität:\\ Ich $=$ Ich, Selbstbewusstsein\\ \footnotesize S.~96~f.}/{Nichtidentität:\\ Ich und Nicht-Ich, Bewusstsein\\ \footnotesize S.~50}/{Fichte: Synthese nur als Sollen},
  6.8/O/{\textbf{Objekt}}/{Identität:\\ Naturzweck, Organisation,\\ \glqq Begriff und Sein identisch\grqq{}\\ \footnotesize S.~103}/{Nichtidentität:\\ Mechanismus, Kausalität\\ \footnotesize S.~38, 103}/{Kant: Zweck nur als Maxime}}
{
  \node[sp] (b\nm) at (\x,0) {};
  \node[above=1mm of b\nm.north, font=\small] {\name};
  \draw[dashed] (b\nm.west) -- (b\nm.east);
  \node[half] at ($(b\nm.center)+(0,0.85)$) {\idt};
  \node[half] at ($(b\nm.center)+(0,-0.85)$) {\nid};
  \node[below=1mm of b\nm.south, font=\scriptsize\itshape] {\fehl};
}
\draw[<->, >={Latex[length=2mm]}, thick, gray!70!black] (bA.east) -- node[above=0.5mm, font=\scriptsize] {gebiert /\\ gebiert sich aus} (bS.west);
\draw[<->, >={Latex[length=2mm]}, thick, gray!70!black] (bS.east) -- node[above=0.5mm, font=\scriptsize] {nur indem das\\ Objekt selbst\\ \SO{} ist} (bO.west);
\end{tikzpicture}}
\caption{Identität und Nichtidentität in den drei Räumen der ersten Rekursion. In jedem Raum wandelt die Unterscheidung ihre Form: Punkt und Ganzes, Selbstbewusstsein und Bewusstsein, Organisation und Mechanismus. Fichte und Kant scheitern spiegelbildlich, weil sie die eine Rekursion ohne die andere führen.}
\label{fig:raeume}
\end{figure}

\subsection{Subjekt und Objekt innerhalb der Identität und der Nichtidentität}

Die Gegenrichtung ist im Text ebenso ausdrücklich, und sie ist noch überraschender, weil sie im Grundsatzkapitel steht, wo von Subjekt und Objekt zunächst gar nicht die Rede ist.

Innerhalb der Identität: Als die Vernunft das Setzen dessen postuliert, wovon in $A = A$ abstrahiert wurde, sagt Hegel: \glqq das eine $A$ ist Subjekt, das andere Objekt\grqq{} \Werke{37}. Die beiden $A$ des Identitätssatzes, die für den Verstand bloße Wiederholung sind, sind für die Vernunft Subjekt und Objekt. Die Unterscheidung von Subjekt und Objekt tritt in den Raum der Identität ein als die Differenz der beiden Stellen im Satz, links und rechts vom Gleichheitszeichen. Wer $A = A$ nur als Verstandesidentität liest, sieht sie nicht; wer den Satz als Beziehung liest, sieht, dass er \glqq die Differenz des $A$ als Subjekts und $A$ als Objekts zugleich mit der Identität\grqq{} enthält \Werke{39}.

Innerhalb der Nichtidentität: Der Satz $A = B$ setzt die Nichtidentität, aber er kann sie nur setzen, weil das Nichtdenken gedacht wird; das Gleichheitszeichen ist in ihm \glqq nur subjektiv\grqq{}, und \glqq dies Gesetztsein des Nichtdenkens fürs Denken ist dem Nichtdenken durchaus zufällig\grqq{} \Werke{38}. Hier tritt die Unterscheidung von Subjekt und Objekt in den Raum der Nichtidentität ein als Denken und Nichtdenken: Das Denken, das die Differenz setzt, ist die subjektive Seite; das Nichtdenken, dem dieses Gesetztsein zufällig ist, die objektive. Die Nichtidentität ist innerlich gespalten in das Setzen der Differenz und die Differenz, die dem Setzen gleichgültig ist.

Innerhalb der Antinomie schließlich, der Beziehung beider Sätze, ist die Unterscheidung von Subjekt und Objekt als die Richtung des Bezogenseins da: Der zweite Satz auf den ersten bezogen sagt, der Widerspruch sei so notwendig wie der Nichtwiderspruch; der erste auf den zweiten bezogen enthält die Differenz mit der Identität. Beide Richtungen zusammen sind die Antinomie, und als Antinomie ist es gleichgültig, welcher Satz gesetzt wird \Werke{39}. Das ist die Indifferenz, in der Form des Erkennens.

\subsection{Die Einheit}

Nun lässt sich sagen, worin die Einheit der beiden Unterscheidungen besteht. Sie sind nicht zwei Rekursionen, die nebeneinander laufen, sondern eine Rekursion in zwei Formen, und die Formen sind die von Hegel benannten: Sein und Erkennen. Die Unterscheidung von Subjekt und Objekt ist die Unterscheidung, wie sie ist: als Natur und Intelligenz, als die beiden Wissenschaften, als Pole der Indifferenz. Die Unterscheidung von Identität und Nichtidentität ist dieselbe Unterscheidung, wie sie gewusst wird: als die beiden Sätze, als Antinomie, als Anschauung. Was beide zu einer macht, ist die Vernunft, die \glqq Identität und Nichtidentität identisch setzt\grqq{} in beiden Formen zugleich \Werke{98}, und die das kann, weil sie es selbst ist, die sich als Natur und als Intelligenz produziert und sich in ihnen erkennt \Werke{100}.

Der Beweis dafür ist die wechselseitige Übersetzbarkeit, die sich gezeigt hat. Jede Unterscheidung tritt in die Räume der anderen ein, ohne Rest, aber nicht ohne Wandlung der Form. Identität heißt im Absoluten Punkt, im Subjekt Selbstbewusstsein, im Objekt Organisation; Nichtidentität heißt dort Ganzes, Bewusstsein, Mechanismus. Subjekt und Objekt heißen in der Identität die beiden Stellen des Satzes, in der Nichtidentität Denken und Nichtdenken, in der Antinomie die beiden Richtungen des Bezogenseins. Wären die beiden Unterscheidungen zwei, so müsste eine von ihnen in einem der Räume der anderen fehlen oder ihre Form dort unverändert behalten. Beides ist nicht der Fall. Sie fehlen nirgends, und sie behalten ihre Form nirgends.

Die beiden Unterschiede zwischen den Rekursionen, die im zweiten Teil festgestellt wurden, klären sich von hier aus. Die Richtung: Die erste Rekursion ist Darstellung von oben, die zweite Synthese von unten. Das sind die beiden Richtungen des einen Satzes, das Absolute gebäre beide und gebäre sich aus beiden \Werke{94}. Es gebiert beide, indem es sich als Subjekt und Objekt darstellt; es gebiert sich aus beiden, indem es als Identität der beiden Sätze hervorgeht. Die Rekursion des Seins läuft abwärts, die des Erkennens aufwärts, und beide sind ein Organismus, dessen Ganzes die Teile und dessen Teile das Ganze hervorbringen. Die Namensgebung: Die Wurzel des Seins heißt \SO{}, mit beiden Namen, weil im Sein beide Seiten in gleicher Würde bestehen. Die Wurzel des Erkennens heißt Identität, mit einem Namen, weil das Erkennen, das die beiden Sätze zusammenfasst, die Anschauung ist, und die Anschauung das Zugleich als Eines hat. Die Asymmetrie des Namens ist die Asymmetrie zwischen dem, was ist, und dem, was es weiß.

Und damit ist auch die Frage nach dem Und/Oder beantwortet. An der Wurzel steht in beiden Formen das Und: Subjekt und Objekt, Identität und Nichtidentität, zugleich. Das Oder gehört dem Verstand, der einen der beiden Sätze gelten lassen will, und Fichte und Kant sind die beiden Gestalten dieses Oder: Fichte lässt im Subjekt nur die Identität gelten und macht die Nichtidentität zum Sollen; Kant lässt im Objekt nur die Nichtidentität gelten und macht die Identität zur Maxime. Das Und der Vernunft ist die Identität der Identität und der Nichtidentität, und diese ist, in der Form des Seins gesagt, das \SO{}.

%=====================================================================
\section{Ergebnis}

Die Differenzschrift enthält zwei Rekursionen derselben Gestalt. Die erste: Das Absolute ist \SO{}, Subjekt und Objekt sind je \SO{} unter ihrer Form, und die Teilung geht ins Unendliche. Die zweite: Das Absolute ist die Identität der Identität und der Nichtidentität, und die beiden Sätze $A = A$ und $A = B$ enthalten je beide Momente unter ihrer Form und bedingen einander. Beide Rekursionen sind eine, in der Form des Seins und in der Form des Erkennens. Der Operator an der Wurzel ist in beiden das Und, das Hegel mit dem Namen der Identität benennt, weil die Anschauung das Zugleich als Eines hat. Jede der beiden Unterscheidungen tritt in die Räume der anderen ein und wandelt dort ihre Form: Identität und Nichtidentität werden im Absoluten Punkt und Ganzes, im Subjekt Selbstbewusstsein und Bewusstsein, im Objekt Organisation und Mechanismus; Subjekt und Objekt werden in der Identität die beiden Stellen des Satzes und in der Nichtidentität Denken und Nichtdenken. Fichte und Kant sind die beiden Weisen, eine Rekursion ohne die andere zu führen, und ihre Ergebnisse, das Sollen und die Maxime, sind die beiden Gestalten der nicht vollzogenen Einheit. Der Satz, in dem beide Rekursionen sich schneiden, lautet: Nur indem das Objekt selbst ein \SO{} ist, ist Ich~=~Ich das Absolute.

\clearpage
\section*{Apparat}

\subsection*{1. Statustafel}
\begin{center}
\renewcommand{\arraystretch}{1.3}
\begin{tabular}{@{}R{8.6cm}R{2.6cm}R{3.4cm}@{}}
\toprule
Aussage & Status & Beleg \\
\midrule
Subjekt und Objekt werden beide als \SO{} gesetzt; die Teilung geht ins Unendliche. & belegt & W~2, S.~94, 96~f. \\
Der Satz der Identität enthält die Differenz, der Satz der Nichtidentität die Identität; beide bedingen einander. & belegt & W~2, S.~37--39 \\
Die zweite Rekursion hat dieselbe Gestalt wie die erste (Paar als Wurzel, Kinder je das Paar unter einer Form). & Deutung, gestützt & W~2, S.~37--39, 96 \\
Alle Kombinationen von Wurzeloperator und Argument sind im Text benannt. & belegt & W~2, S.~20, 37--39, 94, 96, 97 \\
Die beiden Spalten verhalten sich wie Form des Seins zu Form des Erkennens. & Deutung, gestützt & W~2, S.~98 \\
Identität/Nichtidentität im Absoluten $=$ Identität/Totalität. & Deutung, gestützt & W~2, S.~99, 111 \\
Identität/Nichtidentität im Subjekt $=$ Ich~=~Ich / Ich und Nicht-Ich. & belegt & W~2, S.~50, 96~f. \\
Identität/Nichtidentität im Objekt $=$ Naturzweck / Mechanismus. & belegt & W~2, S.~38, 103 \\
Subjekt/Objekt in der Identität $=$ die beiden $A$; in der Nichtidentität $=$ Denken/Nichtdenken. & belegt & W~2, S.~37, 38 \\
Die Asymmetrie der Wurzelnamen ist die von Sein und Erkennen. & Deutung & W~2, S.~94, 96, 98 \\
Fichte und Kant als spiegelbildliche Gestalten des Oder. & Deutung, gestützt & W~2, S.~50, 103 \\
Die offene Strecke bildet den gedoppelten Mittelpunkt ab. & Deutung, rechnerisch geprüft & W~2, S.~111 \\
Seitenangabe zur Bedingtheit des Ich durch das Objekt. & unscharf (S.~68--73) & Paginierung der Vorlage lückenhaft \\
\bottomrule
\end{tabular}
\end{center}

\subsection*{2. Abbildungen}
\begin{itemize}
\item Abbildung~\ref{fig:baum}: Die erste Rekursion als Baum mit doppelter Erzeugungsrichtung.
\item Abbildung~\ref{fig:strecke}: Die offene Strecke und ihre Teilungen; der gedoppelte Mittelpunkt.
\item Abbildung~\ref{fig:baumIN}: Die zweite Rekursion als Baum derselben Gestalt.
\item Abbildung~\ref{fig:matrix}: Wurzeloperator und Argument; die Kombinationen mit ihren Namen im Text.
\item Abbildung~\ref{fig:raeume}: Identität und Nichtidentität in den Räumen des Absoluten, des Subjekts und des Objekts.
\end{itemize}

\subsection*{3. Sachen}
\begin{itemize}
\item \SO{}
\item subjektives \SO{}, objektives \SO{}
\item Indifferenzpunkt
\item Identität und Totalität (gedoppelter Mittelpunkt)
\item Entzweiung
\item ideelle und reelle Entgegensetzung
\item Substantialitätsverhältnis
\item Teilung ins Unendliche
\item offene Strecke, Pol und Mitte
\item Satz der Identität ($A = A$), Satz des Grundes ($A = B$)
\item Antinomie
\item transzendentale Anschauung
\item Verstandesidentität und Vernunftidentität
\item Form des Seins und Form des Erkennens
\item Selbstbewusstsein und Bewusstsein
\item Naturzweck, Organisation, Mechanismus
\item Denken und Nichtdenken
\item Sollen und Maxime
\item Transzendentalphilosophie, Naturphilosophie
\end{itemize}

\subsection*{4. Personen}
\begin{itemize}
\item Georg Wilhelm Friedrich Hegel
\item Johann Gottlieb Fichte
\item Friedrich Wilhelm Joseph Schelling
\item Immanuel Kant
\item Karl Leonhard Reinhold
\end{itemize}

\subsection*{5. Quellen}
\begin{itemize}
\item Hegel, G.\,W.\,F.: Differenz des Fichteschen und Schellingschen Systems der Philosophie (1801). In: Werke, Band~2: Jenaer Schriften 1801--1807. Frankfurt am Main: Suhrkamp 1986, S.~9--138. Zitiert als W~2.
\item Kant, I.: Kritik der Urteilskraft (1790), §§~64--65, zur Bestimmung organisierter Wesen als Naturzwecke. Sinngemäß wiedergegeben.
\item Fichte, J.\,G.: Grundlage der gesamten Wissenschaftslehre. Sämtliche Werke, Band~I. Nach Hegels Darstellung, W~2, S.~52~ff.
\end{itemize}

\subsection*{Dokumentdaten}
\begin{itemize}
\item Titel: Zwei Rekursionen. Subjekt-Objekt und Identität und Nichtidentität in Hegels Differenzschrift (1801) und die Einheit beider Unterscheidungen
\item Datei: zwei\_rekursionen\_subjekt-objekt\_und\_identitaet-nichtidentitaet\_v1.tex
\item Version: v1
\item Datum: 2026-09-25
\item Herkunft: zweite Fassung des Essays \glqq Subjekt-Objekt und die Identität der Identität und der Nichtidentität\grqq{} (v2); der erste und zweite Teil übernehmen dessen Befunde unter dem Gesichtspunkt der Rekursion, der dritte Teil ist neu
\item Grundlage: Hegel, Werke Band~2 (Suhrkamp 1986), Textdatei der Vorlage; Seitenzahlen nach eingedruckter Paginierung
\item Status: Erstfassung; Kanon nach Pauschale, solange die Qualität hält
\end{itemize}

\end{document}
